\chapter{1995年上海卷（理）}
\section{单选题}

\begin{problem}
\(y=\sin^2 x\) 是（\quad）

\choices
    {最小正周期为 \(2\pi\) 的偶函数}
    {最小正周期为 \(2\pi\) 的奇函数}
    {最小正周期为 \(\pi\) 的偶函数}
    {最小正周期为 \(\pi\) 的奇函数}
\end{problem}

\begin{answer}
C
\end{answer}

\begin{problem}
如果 \(P=\{x\mid (x-1)(2x-5)<0\}\)，\(Q=\{x\mid 0<x<10\}\)，那么（\quad）

\choices
    {\(P\cap Q=\emptyset\)}
    {\(P\subset Q\)}
    {\(P\supset Q\)}
    {\(P\cup Q=\R\)}
\end{problem}

\begin{answer}
B
\end{answer}

\begin{problem}
方程 \(\tan\left(2x+\frac{\pi}{3}\right)=\frac{\sqrt{3}}{3}\) 在区间 \([0,2\pi)\) 上解的个数是（\quad）

\choices
    {5}
    {4}
    {3}
    {2}
\end{problem}

\begin{answer}
B
\end{answer}

\begin{problem}
设棱锥的底面面积是 \(8\mathrm{~cm}^2\)，那么这个棱锥的中截面（过棱锥高的中点且平行于底面的截面）的面积是（\quad）

\choices
    {\(4\mathrm{~cm}^2\)}
    {\(2\sqrt{2}\mathrm{~cm}^2\)}
    {\(2\mathrm{~cm}^2\)}
    {\(\sqrt{2}\mathrm{~cm}^2\)}
\end{problem}

\begin{answer}
C
\end{answer}

\begin{problem}
“\(ab<0\)”是“方程 \(ax^2+by^2=c\) 表示双曲线”的（\quad）

\choices
    {必要条件但不是充分条件}
    {充分条件但不是必要条件}
    {充分必要条件}
    {既不是充分条件又不是必要条件}
\end{problem}

\begin{answer}
A
\end{answer}

\begin{problem}
当 \(a\neq 0\) 时，函数 \(y=ax+b\) 和 \(y=b^{ax}\) 的图象只可能是（\quad）

\choices
    {\choicebitmap[width=0.15\paperwidth]{img/1995/shanghai_science/q6_optA_nolabel.png}}
    {\choicebitmap[width=0.15\paperwidth]{img/1995/shanghai_science/q6_optB_nolabel.png}}
    {\choicebitmap[width=0.15\paperwidth]{img/1995/shanghai_science/q6_optC_nolabel.png}}
    {\choicebitmap[width=0.15\paperwidth]{img/1995/shanghai_science/q6_optD_nolabel.png}}
\end{problem}

\begin{answer}
A
\end{answer}

\begin{problem}
当 \(0<a<b<1\) 时，下列不等式中正确的是（\quad）

\choices
    {\((1-a)^{\frac{1}{b}}>(1-a)^b\)}
    {\((1+a)^a>(1+b)^b\)}
    {\((1-a)^b>(1-a)^{\frac{b}{2}}\)}
    {\((1-a)^a>(1-b)^b\)}
\end{problem}

\begin{answer}
D
\end{answer}

\begin{problem}
下列四个命题中的真命题是（\quad）

\choices
    {经过定点 \(P_0(x_0,y_0)\) 的直线都可以用方程 \(y-y_0=k(x-x_0)\) 表示}
    {经过任意两个不同的点 \(P_1(x_1,y_1)\)，\(P_2(x_2,y_2)\) 的直线都可以用方程 \((y-y_1)(x_2-x_1)=(x-x_1)(y_2-y_1)\) 表示}
    {不经过原点的直线都可以用方程 \(\frac{x}{a}+\frac{y}{b}=1\) 表示}
    {经过定点 \(A(0,b)\) 的直线都可以用方程 \(y=kx+b\) 表示}
\end{problem}

\begin{answer}
B
\end{answer}

\section{填空题}

\begin{problem}
不等式 \(\frac{2x-1}{x+3}>1\) 的解是\fillinblank{}.
\end{problem}

\begin{answer}
\(x<-3\) 或 \(x>4\)
\end{answer}

\begin{problem}
双曲线 \(\frac{x^2}{2}-\frac{8y^2}{9}=8\) 的渐近线方程是\fillinblank{}.
\end{problem}

\begin{answer}
\(3x\pm4y=0\)
\end{answer}

\begin{problem}
1992年底世界人口达到 \(54.8\text{ 亿}\)，若人口的年平均增长率为 \(x\%\)，2000 年底世界人口数为 \(y\)（亿），那么 \(y\) 与 \(x\) 的函数关系式是\fillinblank{}.
\end{problem}

\begin{answer}
\(y=54.8(1+x\%)^8\)
\end{answer}

\begin{problem}
已知 \(\log_3 x=\frac{-1}{\log_2 3}\)，那么 \(x+x^2+x^3+\cdots+x^n+\cdots=\)\fillinblank{}.
\end{problem}

\begin{answer}
\(1\)
\end{answer}

\begin{problem}
若 \((x+1)^n=x^n+\cdots+ax^3+bx^2+\cdots+1\)（\(n\in\N\)），且 \(a:b=3:1\)，那么 \(n=\)\fillinblank{}.
\end{problem}

\begin{answer}
\(11\)
\end{answer}

\begin{problem}
到点 \(A(-1,0)\) 和直线 \(x=3\) 距离相等的点的轨迹方程是\fillinblank{}.
\end{problem}

\begin{answer}
\(y^2=-8x+8\)
\end{answer}

\begin{problem}
把参数方程 \(x=\sin\alpha\)，\(y=\cos\alpha+1\)（\(\alpha\) 是参数）化为普通方程，结果是\fillinblank{}.
\end{problem}

\begin{answer}
\(x^2+(y-1)^2=1\)
\end{answer}

\begin{problem}
把直角坐标系的原点作为极点，\(x\) 轴的正半轴作为极轴，并且在两种坐标系中取相同的长度单位. 若曲线的极坐标方程是 \(\rho^2=\frac{1}{4\cos^2\theta-1}\)，则它的直角坐标方程是\fillinblank{}.
\end{problem}

\begin{answer}
\(3x^2-y^2=1\)
\end{answer}

\begin{problem}
函数 \(y=\sin\frac{x}{2}+\cos\frac{x}{2}\) 在 \((-2\pi,2\pi)\) 内的递增区间是\fillinblank{}.
\end{problem}

\begin{answer}
\(\left[-\frac{3\pi}{2},\frac{\pi}{2}\right]\)
\end{answer}

\section{选做题：填空题}

\begin{problem}[A]
若 \(\displaystyle\lim_{n\to\infty}\left[1+(r+1)^n\right]=1\)，则 \(r\) 的取值范围是\fillinblank{}.
\end{problem}

\begin{answer}
\(-2<r<0\)
\end{answer}

\begin{problem}[A]
从 6 台原装计算机和 5 台组装计算机中任意选取 5 台，其中至少有原装与组装计算机各 2 台，则不同的选取法有\fillinblank{} 种（结果用数值表示）.
\end{problem}

\begin{answer}
\(350\)
\end{answer}

\begin{problem}[A]
把圆心角为 \(216^\circ\)，半径为 \(5\text{ 分米}\) 的扇形铁皮焊成一个圆锥形容器（不计焊缝），那么容器的容积是\fillinblank{}\(\text{立方分米}\)（结果保留两位小数）.
\end{problem}

\begin{answer}
\(37.70\)（或 \(37.68\)，\(37.69\)）
\end{answer}

\setcounter{probnum}{17}
\begin{problem}[B]
\(\displaystyle\lim_{n\to\infty}\left(1+\frac{1}{n}\right)^{n-2}=\)\fillinblank{}.
\end{problem}

\begin{answer}
\(\e\)
\end{answer}

\begin{problem}[B]
从 6 台原装计算机和 5 台组装计算机中任意选取 5 台，其中至少有原装与组装计算机各 2 台的概率是\fillinblank{}.
\end{problem}

\begin{answer}
\(\frac{25}{33}\)（或 \(0.76\)）
\end{answer}

\begin{problem}[B]
复数 \(z\) 满足 \((1+2\i)\overline{z}=4+3\i\)，那么 \(z=\)\fillinblank{}.
\end{problem}

\begin{answer}
\(2+\i\)
\end{answer}

\section{选做题：解答题}

\begin{problem}[A]
已知 \(\tan\left(\frac{\pi}{4}+\theta\right)=3\)，求 \(\sin2\theta-2\cos^2\theta\) 的值.
\end{problem}

\begin{solution}
因为 \(\tan\left(\frac{\pi}{4}+\theta\right)=3\)，所以
\[
\frac{1+\tan\theta}{1-\tan\theta}=3
\]
解得 \(\tan\theta=\frac{1}{2}\).

于是
\[
\sin2\theta-2\cos^2\theta=\sin2\theta-\cos2\theta-1=\frac{2\tan\theta}{1+\tan^2\theta}-\frac{1-\tan^2\theta}{1+\tan^2\theta}-1=\frac{4}{5}-\frac{3}{5}-1=-\frac{4}{5}
\]
\end{solution}

\begin{problem}[A]
已知复数 \(z_1\)，\(z_2\) 满足 \(|z_1|=|z_2|=1\)，且 \(z_1+z_2=\frac{1}{2}+\frac{\sqrt{3}}{2}\i\). 求 \(z_1\)，\(z_2\) 的值.
\end{problem}

\begin{solution}
由题意可设 \(z_1=\cos\alpha+\i\sin\alpha\)，\(z_2=\cos\beta+\i\sin\beta\)，其中 \(\alpha\geqslant\beta\)，且 \(\alpha\)，\(\beta\in[0,2\pi)\).

由 \(z_1+z_2=\frac{1}{2}+\frac{\sqrt{3}}{2}\i\)，得
\[
\begin{cases}
\cos\alpha+\cos\beta=\frac{1}{2},\\
\sin\alpha+\sin\beta=\frac{\sqrt{3}}{2}
\end{cases}
\]

由上式两式平方相加，得 \(\cos(\alpha-\beta)=-\frac{1}{2}\).

又 \(\alpha-\beta\in[0,2\pi)\)，所以 \(\alpha-\beta=\frac{2\pi}{3}\) 或 \(\alpha-\beta=\frac{4\pi}{3}\).

当 \(\alpha-\beta=\frac{2\pi}{3}\) 时，把 \(\alpha=\beta+\frac{2\pi}{3}\) 代入上式，得 \(\cos\beta=1\)，\(\sin\beta=0\)，从而 \(z_1=-\frac{1}{2}+\frac{\sqrt{3}}{2}\i\)，\(z_2=1\).

当 \(\alpha-\beta=\frac{4\pi}{3}\) 时，把 \(\alpha=\beta+\frac{4\pi}{3}\) 代入上式，得 \(\cos\beta=-\frac{1}{2}\)，\(\sin\beta=\frac{\sqrt{3}}{2}\)，从而 \(z_1=1\)，\(z_2=-\frac{1}{2}+\frac{\sqrt{3}}{2}\i\).

因此
\[
\begin{cases}
z_1=1,\\
z_2=-\frac{1}{2}+\frac{\sqrt{3}}{2}\i
\end{cases}
\]

或

\[
\begin{cases}
z_1=-\frac{1}{2}+\frac{\sqrt{3}}{2}\i,\\
z_2=1.
\end{cases}
\]

解二：由 \(|z_1+z_2|=1\)，得
\[
(z_1+z_2)(\overline{z_1}+\overline{z_2})=1
\]
又 \(|z_1|=|z_2|=1\)，故
\[
\overline{z_1}z_2+\overline{z_2}z_1=-1
\]
所以 \(\overline{z_1}z_2\) 的实部和 \(\overline{z_2}z_1\) 的实部均为 \(-\frac{1}{2}\).

又 \(|\overline{z_1}z_2|=1\)，故 \(\overline{z_1}z_2\) 的虚部为 \(\pm\frac{\sqrt{3}}{2}\)，即
\[
\overline{z_1}z_2=-\frac{1}{2}\pm\frac{\sqrt{3}}{2}\i
\]
于是
\[
z_2=z_1\overline{z_1}z_2=z_1\left(-\frac{1}{2}\pm\frac{\sqrt{3}}{2}\i\right)
\]
所以
\[
z_1+z_2\left(-\frac{1}{2}\pm\frac{\sqrt{3}}{2}\i\right)=\frac{1}{2}+\frac{\sqrt{3}}{2}\i
\]
从而
\[
z_1=1,\quad z_2=-\frac{1}{2}+\frac{\sqrt{3}}{2}\i
\]
或
\[
z_1=-\frac{1}{2}+\frac{\sqrt{3}}{2}\i,\quad z_2=1.
\]
\end{solution}

\setcounter{probnum}{20}
\begin{problem}[B]
如图，在空间直角坐标系中，\(BC=2\)，原点 \(O\) 是 \(BC\) 的中点，点 \(A\) 的坐标是 \(\left(\frac{\sqrt{3}}{2},\frac{1}{2},0\right)\)，点 \(D\) 在平面 \(yoz\) 上，且 \(\angle BDC=90^\circ\)，\(\angle DCB=30^\circ\).

\begin{enumerate}
\item 求向量 \(\overrightarrow{OD}\) 的坐标；
\item 设向量 \(\overrightarrow{AD}\) 和 \(\overrightarrow{BC}\) 的夹角为 \(\theta\)，求 \(\cos\theta\) 的值.
\end{enumerate}

\bitmapfigure[max width=0.42\linewidth]{img/1995/shanghai_science/q21_fig1.png}
\end{problem}

\begin{solution}
\begin{enumerate}
\item \(z_D=|BC|\cos30^\circ\sin30^\circ=\frac{\sqrt{3}}{2}\)，\(y_D=1-|BC|\cos30^\circ\cos30^\circ=-\frac{1}{2}\)，所以
\(\overrightarrow{OD}=\left(0,-\frac{1}{2},\frac{\sqrt{3}}{2}\right)\).

\item
由 \(\overrightarrow{OA}=\left(\frac{\sqrt{3}}{2},\frac{1}{2},0\right)\)，
\(\overrightarrow{OB}=(0,-1,0)\)，
\(\overrightarrow{OC}=(0,1,0)\)，得
\(\overrightarrow{AD}=\left(-\frac{\sqrt{3}}{2},-1,\frac{\sqrt{3}}{2}\right)\)，
\(\overrightarrow{BC}=(0,2,0)\).

于是
\[
\cos\theta=\frac{\overrightarrow{AD}\cdot\overrightarrow{BC}}{|\overrightarrow{AD}|\,|\overrightarrow{BC}|}
=\frac{-2}{\sqrt{\frac{5}{2}}\cdot2}=-\frac{\sqrt{10}}{5}
\]
\end{enumerate}
\end{solution}

\begin{problem}[B]
设 \(y=f(x)\) 是二次函数，方程 \(f(x)=0\) 有两个相等的实根，且 \(f'(x)=2x+2\).

\begin{enumerate}
\item 求 \(y=f(x)\) 的表达式；
\item 若直线 \(x=-t\)（\(0<t<1\)）把 \(y=f(x)\) 的图象与两坐标轴所围成图形的面积二等分，求 \(t\) 的值.
\end{enumerate}
\end{problem}

\begin{solution}
\begin{enumerate}
\item 设 \(f(x)=ax^2+bx+c\)，则 \(f'(x)=2ax+b\). 又已知 \(f'(x)=2x+2\)，故 \(a=1\)，\(b=2\).

于是 \(f(x)=x^2+2x+c\). 方程 \(f(x)=0\) 有两个相等实根，判别式 \(\Delta=4-4c=0\)，得 \(c=1\)，所以 \(f(x)=x^2+2x+1\).

\item 按题意有
\[
\int_{-1}^{-t}(x^2+2x+1)\,\mathrm{d}x=\int_{-t}^{0}(x^2+2x+1)\,\mathrm{d}x
\]
即 \(\left(\frac{1}{3}x^3+x^2+x\right)\Big|_{-1}^{-t}=\left(\frac{1}{3}x^3+x^2+x\right)\Big|_{-t}^{0}\).
\[
-\frac{1}{3}t^3+t^2-t+\frac{1}{3}=\frac{1}{3}t^3-t^2+t
\]
\[
2t^3-6t^2+6t-1=0
\]
即 \(2(t-1)^3=-1\)，于是 \(t=1-\frac{1}{\sqrt[3]{2}}\).
\end{enumerate}
\end{solution}

\begin{problem}
如图，四棱锥 \(P-ABCD\) 中，底面是一个矩形，\(AB=3\)，\(AD=1\)，又 \(PA\perp AB\)，\(PA=4\)，\(\angle PAD=60^\circ\).

\begin{enumerate}
\item 求四棱锥 \(P-ABCD\) 的体积；
\item 求二面角 \(P-BC-D\) 的大小（用反三角函数表示）.
\end{enumerate}

\bitmapfigure[max width=0.34\linewidth]{img/1995/shanghai_science/q23_fig1.png}
\end{problem}

\begin{solution}
\begin{enumerate}
\item 因为 \(AB\perp AD\)，\(AB\perp AP\)，所以 \(AB\perp\) 平面 \(PAD\). 又 \(AB\subset\) 平面 \(ABCD\)，所以平面 \(ABCD\perp\) 平面 \(PAD\).

在平面 \(PAD\) 中，作 \(PE\perp AD\)（\(AD\) 为平面 \(PAD\) 与平面 \(ABCD\) 的交线），交 \(AD\) 的延长线于点 \(E\)（因为 \(AE=AP\cos60^\circ=2>AD\)）.

所以 \(PE\perp\) 平面 \(ABCD\). 在直角三角形 \(PAE\) 中，\(PE=AP\sin60^\circ=2\sqrt{3}\)，故
\[
V_{P-ABCD}=\frac{1}{3}AB\cdot AD\cdot PE=2\sqrt{3}
\]

\item 在平面 \(ABCD\) 中，作 \(EF\parallel DC\)，交 \(BC\) 的延长线于点 \(F\)，则 \(EF\perp BF\)，连接 \(PF\).

因为 \(PE\perp\) 平面 \(ABCD\)，\(EF\perp BF\)，所以 \(PF\perp BF\)（三垂线定理）.

于是 \(\angle PFE\) 是二面角 \(P-BC-D\) 的平面角. 在直角三角形 \(PEF\) 中，
\[
\tan\angle PFE=\frac{PE}{EF}=\frac{2\sqrt{3}}{3}
\]
所以 \(\angle PFE=\arctan\frac{2\sqrt{3}}{3}\). 因此，二面角 \(P-BC-D\) 的大小为 \(\arctan\frac{2\sqrt{3}}{3}\).
\end{enumerate}
\end{solution}

\begin{problem}
设椭圆的方程为 \(\frac{x^2}{m^2}+\frac{y^2}{n^2}=1\)（\(m,n>0\)），过原点且倾角为 \(\theta\) 和 \(\pi-\theta\)（\(0<\theta<\frac{\pi}{2}\)）的两条直线分别交椭圆于 \(A\)，\(C\) 和 \(B\)，\(D\) 四点.

\begin{enumerate}
\item 用 \(\theta\)，\(m\)，\(n\) 表示四边形 \(ABCD\) 的面积 \(S\)；
\item 若 \(m\)，\(n\) 为定值，当 \(\theta\) 在 \(\left(0,\frac{\pi}{4}\right]\) 上变化时，求 \(S\) 的最大值 \(u\)；
\item 如果 \(u>mn\)，求 \(\frac{m}{n}\) 的取值范围.
\end{enumerate}
\end{problem}

\begin{solution}
\begin{enumerate}
\item 设经过原点且倾角为 \(\theta\) 的直线方程为 \(y=x\tan\theta\)，解方程组
\[
\begin{cases}
y=x\tan\theta,\\
\dfrac{x^2}{m^2}+\dfrac{y^2}{n^2}=1
\end{cases}
\]
得 \(\begin{cases}
x^2=\dfrac{m^2n^2}{n^2+m^2\tan^2\theta},\\
y^2=\dfrac{m^2n^2\tan^2\theta}{n^2+m^2\tan^2\theta}
\end{cases}\).
由对称性，四边形 \(ABCD\) 为矩形，又因为 \(0<\theta<\frac{\pi}{2}\)，所以
\[
S=4|xy|=\frac{4m^2n^2\tan\theta}{n^2+m^2\tan^2\theta}
\]

\item
\[
S=\frac{4m^2n^2}{\dfrac{n^2}{\tan\theta}+m^2\tan\theta}
\]

当 \(m>n\) 时，\(\frac{n}{m}<1\). 由基本不等式，
\[
\frac{n^2}{\tan\theta}+m^2\tan\theta\geqslant2mn
\]
且当 \(\tan\theta=\frac{n}{m}\) 时等号成立，所以
\[
S=\frac{4m^2n^2}{\dfrac{n^2}{\tan\theta}+m^2\tan\theta}\leqslant\frac{4m^2n^2}{2mn}=2mn.
\]
由于 \(0<\theta\leqslant\frac{\pi}{4}\)，\(0<\tan\theta\leqslant1\)，故当 \(\tan\theta=\frac{n}{m}\) 时等号成立，得 \(u=2mn\).

当 \(m<n\) 时，\(\frac{n}{m}>1\). 对于任意 \(0<\theta_1<\theta_2\leqslant\frac{\pi}{4}\)，由于
\[
\begin{aligned}
&\left(m^2\tan\theta_2+\frac{n^2}{\tan\theta_2}\right)-\left(m^2\tan\theta_1+\frac{n^2}{\tan\theta_1}\right)\\
&=m^2(\tan\theta_2-\tan\theta_1)-n^2\frac{\tan\theta_2-\tan\theta_1}{\tan\theta_1\tan\theta_2}\\
&=(\tan\theta_2-\tan\theta_1)\frac{m^2\tan\theta_1\tan\theta_2-n^2}{\tan\theta_1\tan\theta_2}.
\end{aligned}
\]
因为 \(0<\tan\theta_1<\tan\theta_2\leqslant1\)，且
\[
m^2\tan\theta_1\tan\theta_2-n^2<m^2-n^2<0,
\]
所以
\[
\left(m^2\tan\theta_2+\frac{n^2}{\tan\theta_2}\right)-\left(m^2\tan\theta_1+\frac{n^2}{\tan\theta_1}\right)<0.
\]
于是
\[
\frac{4m^2n^2}{m^2\tan\theta_2+\dfrac{n^2}{\tan\theta_2}}-\frac{4m^2n^2}{m^2\tan\theta_1+\dfrac{n^2}{\tan\theta_1}}>0.
\]
所以在 \(\left(0,\frac{\pi}{4}\right]\) 上，\(S=\frac{4m^2n^2}{\frac{n^2}{\tan\theta}+m^2\tan\theta}\) 是 \(\theta\) 的递增函数，故当 \(\theta=\frac{\pi}{4}\) 时，\(\tan\theta=1\)，得 \(u=\frac{4m^2n^2}{m^2+n^2}\).

综上，
\[
u=\begin{cases}
2mn,&0<n<m,\\
\dfrac{4m^2n^2}{m^2+n^2},&0<m<n.
\end{cases}
\]

\item 当 \(\frac{m}{n}>1\) 时，\(u=2mn>mn\) 恒成立.

当 \(\frac{m}{n}<1\) 时，
\[
\frac{u}{mn}=\frac{4mn}{m^2+n^2}>1
\]
即 \(\left(\frac{m}{n}\right)^2-4\frac{m}{n}+1<0\).
所以 \(2-\sqrt{3}<\frac{m}{n}<2+\sqrt{3}\). 又由 \(\frac{m}{n}<1\)，得 \(2-\sqrt{3}<\frac{m}{n}<1\).

综上所述，当 \(u>mn\) 时，\(\frac{m}{n}\) 的取值范围为
\[
\left(2-\sqrt{3},1\right)\cup\left(1,+\infty\right)
\]
\end{enumerate}
\end{solution}

\begin{problem}
已知二次函数 \(y=f(x)\) 在 \(x=\frac{t+2}{2}\) 处取得最小值 \(-\frac{t^2}{4}\)（\(t>0\)），\(f(1)=0\).

\begin{enumerate}
\item 求 \(y=f(x)\) 的表达式；
\item 若任意实数 \(x\) 都满足等式 \(f(x)g(x)+a_nx+b_n=x^{n+1}\)（\(g(x)\) 为多项式，\(n\in\N\)），试用 \(t\) 表示 \(a_n\) 和 \(b_n\)；
\item 设圆 \(C_n\) 的方程为 \((x-a_n)^2+(y-b_n)^2=r_n^2\)，圆 \(C_n\) 与 \(C_{n+1}\) 外切（\(n=1\)，\(2\)，\(3\)，\(\cdots\)），\(\{r_n\}\) 是各项都是正数的等比数列，记 \(S_n\) 为前 \(n\) 个圆的面积之和，求 \(r_n\)，\(S_n\).
\end{enumerate}
\end{problem}

\begin{solution}
\begin{enumerate}
\item 设
\[
f(x)=a\left(x-\frac{t+2}{2}\right)^2-\frac{t^2}{4}
\]
由 \(f(1)=0\)，得
\[
a\left(1-\frac{t+2}{2}\right)^2-\frac{t^2}{4}=0
\]
得 \(a=1\)，于是
\[
f(x)=x^2-(t+2)x+t+1
\]

\item \(f(x)=(x-1)[x-(t+1)]\)，代入已知等式，得
\[
(x-1)[x-(t+1)]g(x)+a_nx+b_n=x^{n+1}
\]
将 \(x=1\)，\(x=t+1\) 分别代入上式，得
\[
\begin{cases}
a_n+b_n=1,\\
(t+1)a_n+b_n=(t+1)^{n+1}
\end{cases}
\]
由 \(t\neq0\)，可解得
\[
a_n=\frac{1}{t}\left[(t+1)^{n+1}-1\right],\qquad b_n=\frac{t+1}{t}\left[1-(t+1)^n\right]
\]

\item \(a_n+b_n=1\)，故圆 \(C_n\) 的圆心 \(O_n\) 在直线 \(x+y=1\) 上，于是
\[
|O_nO_{n+1}|=\sqrt{2}|a_{n+1}-a_n|
\]
又圆 \(C_n\) 与 \(C_{n+1}\) 外切，故
\[
r_n+r_{n+1}=\sqrt{2}(t+1)^{n+1}
\]
设 \(\{r_n\}\) 的公比为 \(q\)，则
\[
\begin{cases}
r_n+qr_n=\sqrt{2}(t+1)^{n+1},\\
r_{n+1}+qr_{n+1}=\sqrt{2}(t+1)^{n+2}
\end{cases}
\]
两式相除，得 \(q=\frac{r_{n+1}}{r_n}=t+1\).

于是
\[
r_n=\frac{\sqrt{2}(t+1)^{n+1}}{t+2}
\]
\[
S_n=\pi(r_1^2+r_2^2+\cdots+r_n^2)=\frac{\pi r_1^2(q^{2n}-1)}{q^2-1}=\frac{2\pi(t+1)^4}{t(t+2)^3}\left[(t+1)^{2n}-1\right]
\]
\end{enumerate}
\end{solution}
